% SPDX-License-Identifier: GPL-3.0-or-later
% French -- what the reading editions' preamble leaves to the language.
% Input by lib/frank-preamble.tex as \input{\FrankLib/lang/\BookLang.tex};
% docs/languages.md section 6 lists the hooks every file here must define.
% Latin script, like Italian: no font of its own, nothing to strip, two
% passes.  What is not like Italian is the how-to page.  French spelling
% records a pronunciation the language has moved away from, so the middle
% line of every gloss is doing real work, and the reader has to be told how
% to read it before the first chapter rather than at the first surprise.

% ---- the babel language and the switches ----------------------------------
\newcommand{\FrankLangName}{french}
\FrankRTLfalse           % chunk left, gloss right; \pw needs no bidi isolate
\FrankHasBarefalse       % there is nothing to strip, so pass 3 would repeat pass 1
\FrankHasAltfalse        % no alternate face
\FrankHasReadingfalse
\FrankHasVertfalse

% ---- fonts ------------------------------------------------------------------
% The registry's tex.main is null for French: the text is set in the main
% roman the core preamble already loaded (TeX Gyre Pagella), the same face as
% the glosses.  Pagella carries the whole of written French -- the accents,
% the cedilla, the oe ligature -- and the letters the pronunciation line adds
% to it (ã ẽ õ ò ü ʃ ʒ ɲ) come from the DejaVu Serif fallback the core declares
% glyph by glyph; checked on this machine: the fixture edition builds with no
% Missing character in its log.
% import=fr asks babel for the French locale, which brings the hyphenation
% patterns of hyph-fr when the machine has them.  install.sh --pdf installs
% them (hyphen-french) and says so; what a machine without them does is leave
% \language at 0 and break a French paragraph at the English patterns' points
% -- print-emps for prin-temps, répon-dit for ré-pon-dit -- silently, so it is
% install.sh and build.sh that have to say it, and both now do.
%% transforms=punctuation.space is what makes the spacing French.  Without
%% it \babelprovide loads the fr locale but not french.ldf, so `Oui?' and
%% `Oui ?' set 23.2pt and 26.0pt -- English spacing, or a full word space --
%% and the guillemets sit tight or loose according to what was typed.  With
%% it both come out at 24.6pt, a thin space, and `<<mot>>' and `<< mot >>'
%% both at 34.4pt: whatever the source did, the page is right.  That is what
%% lets docs/lang/fr.md tell the annotator to reproduce the source and
%% nothing more.  Measured with lualatex at 11pt in TeX Gyre Pagella.
\babelprovide[import=fr,transforms=punctuation.space]{french}

% ---- the vocabulary line -----------------------------------------------------
% \vb{infinitive}{sound}{first person present}{sound}{past participle}{sound}{meaning}
% The three forms a French verb cannot be rebuilt from any one of.  The present
% is where the stem breaks (aller but je vais, pouvoir but je peux, boire but
% je bois); the participle is what every compound past is made of (j'ai fait,
% elle est venue) and what agrees after être.  The imperfect, the future and
% the conditional follow a pattern from those two, so they get no slot -- and
% the past historic, which follows nothing, is named in the gloss where a book
% uses it (docs/lang/fr.md, the vocabulary section).
% What the three cannot say goes in ONE parenthesis after the meaning, items
% parted by "; ", each only where it is not the ordinary case:
%   aux. être, or aux. être/avoir     avoir is the default and never printed
%   nous \pw{prenons}                 where the plural is not the infinitive's stem + -ons
%   fut. \pw{irai}                    where the future is not the infinitive's
%   \vb{aller}{alé}{vais}{vè}{allé}{alé}{to go (aux. être; fut. \pw{irai} \textit{iré})}
% The class the gloss once carried (-er, -ir, -re) is not among them: the
% infinitive prints its own ending, and it was wrong for aller anyway.
% The reader takes the same pair from the registry (lib/languages.json,
% "vb_labels": ["pres.", "p.p."], the three forms in words as "vb_forms"); the
% two must agree, or the PDF and the reader of one book label the same form
% differently.  lib/newlang.py --check compares them.
\newcommand{\FrankVbPres}{pres.}
\newcommand{\FrankVbPast}{p.p.}

% ---- the two Lua filters ---------------------------------------------------
% Both are the identity.  A French accent belongs to the spelling and is not
% help written in for pass 1 -- ou and où, a and à, sur and sûr are different
% words -- so there is nothing a bare pass could take off without changing what
% the page says (the registry's `strip` is null), and the labels are already in
% Latin digits.  They still have to exist, because the core calls them for
% every language.
\directlua{
  function frank_strip(s) return s end
  function frank_latin(s) return s end
}

% ---- the "How to read this" page -----------------------------------------------
% Printed by lib/frank-frontmatter.tex under its heading.  \pw sets a word as
% French inside the English and \textit sets a sound, which is the division
% the gloss line itself makes (\dw is \pw then \textit): so the reader meets
% the two typographies here in the same roles they have on every later page.
% \pw is defined by the core preamble, after this file is read, and is only
% expanded when the page is set.
\newcommand{\FrankHowTo}{%
Every passage comes twice.

\smallskip
\textbf{First} the whole sentence, and nothing else. Try to read it. Get as far
as you can and do not worry about what you cannot get — the point is the
attempt, made before any help arrives.

\smallskip
\textbf{Second} the same sentence in small chunks, French on the left, and on
the right three lines: how it sounds, then how to look each word up, then what
it means. Now you find out how much you had right. There is no third pass: an
accent belongs to the spelling and not to you (\pw{ou} and \pw{où} are two
words), so there is nothing to take away and a bare pass would reprint the
sentence you have just read.

\smallskip
The first of those lines is how the chunk is said, which in French is far from
how it is written. Silent letters are not there: the \textit{t} of \pw{petit},
the \textit{-ent} of \pw{ils parlent}. A consonant silent alone but sounded
before a vowel is carried on a hyphen onto the word it moves to: \pw{les amis}
is \textit{lé-zami}. A tilde is a vowel said through the nose, with no
\textit{n} after it: \textit{ã} in \pw{grand}, \textit{ẽ} in \pw{vin},
\textit{õ} in \pw{bon} — and \pw{bonne} is \textit{bòn}, because there the
\textit{n} really is said. \textit{é} and \textit{è} are the closed and open
\textit{e} of \pw{été} and \pw{mère}, \textit{o} and \textit{ò} the same pair
for \textit{o}, a plain \textit{e} the light vowel of \pw{le}, often
swallowed, and \textit{ü} the \textit{u} of \pw{tu}; \textit{ʃ}, \textit{ʒ}
and \textit{ɲ} are what \textit{ch}, \textit{j} and \textit{gn} spell. It is
left out only where nothing in the chunk can be got wrong.

\smallskip
Verbs are given as infinitive, first person present, past participle, meaning,
each with its sound: the present is where the stem breaks (\pw{aller} makes
\pw{je vais}) and the participle is what every compound past is built on
(\pw{j'ai fait}, \pw{elle est venue}). In brackets after the meaning comes
whatever else the verb will not tell you, and only where it is there to tell:
\textit{aux. être} for a verb whose perfect is made with \pw{être}, and the
\pw{nous} form (\pw{nous prenons}) and the future (\pw{j'irai}) where they
are not built on the infinitive — which is how to build them where they are
not given. A verb welded to a bare noun (\pw{avoir peur}) is given no meaning
of its own; the noun carries the sense.
}
